\documentclass[12pt, a4paper]{article}

\usepackage[utf8]{inputenc}
\usepackage[T1]{fontenc}
\usepackage[german]{babel}
\usepackage{amsmath, amssymb, amsthm}
\usepackage{geometry}
\usepackage{tikz}
\usepackage{tikz-3dplot}
\usepackage{pgfplots}
\usepackage{float}
\usepackage{caption}
\usepackage{subcaption}
\usepackage{booktabs}
\usepackage{array}
\usepackage{hyperref}
\usepackage{xcolor}
\usepackage{mathtools}
\usepackage{graphicx}
\usepackage{rotating}
\usepackage{microtype}
\usepackage{enumitem}
\setlist[itemize,1]{label=--}


\hypersetup{colorlinks=true,linkcolor=blue!60!black,urlcolor=blue!60!black,citecolor=blue!60!black}

\usetikzlibrary{arrows.meta, decorations.pathmorphing, patterns, calc, shapes.geometric,
                shapes.misc, positioning, fit, backgrounds, shadows, fadings}

\pgfplotsset{compat=1.18}

\geometry{left=2.5cm, right=2.5cm, top=3cm, bottom=3cm}

\definecolor{pincolor}{RGB}{70, 130, 180}
\definecolor{scarcolor}{RGB}{220, 80, 60}
\definecolor{ringcolor}{RGB}{100, 160, 100}
\definecolor{lockcolor}{RGB}{255, 200, 50}
\definecolor{unlockcolor}{RGB}{200, 200, 200}
\definecolor{presscolor}{RGB}{180, 100, 220}

% Theorem-Umgebungen
\newtheorem{theorem}{Satz}[section]
\newtheorem{lemma}[theorem]{Lemma}
\newtheorem{corollary}[theorem]{Korollar}
\newtheorem{definition}{Definition}[section]
\newtheorem{proposition}[theorem]{Proposition}

\theoremstyle{remark}
\newtheorem{remark}{Bemerkung}[section]

\title{\textbf{Optimaler Doppelter Ringverschluss}\\
\large Formale Beschreibung und mechanischer Nachweis  eines Zwei-Ebenen-Rastsystems mit 9-Pin-Konfiguration}

\author{Stephan Epp}
\date{\today}

% Hilfsmakros für Formeln
\newcommand{\Pin}{\mathcal{P}}
\newcommand{\Scar}{\mathcal{S}}
\newcommand{\Ring}{\mathcal{R}}
\newcommand{\Lock}{\mathcal{L}}
\newcommand{\Phiangle}{\varphi}
\newcommand{\vvec}[1]{\boldsymbol{#1}}

\begin{document}

\maketitle
\tableofcontents
\newpage

% ============================================================
\section{Einleitung}
% ============================================================

Mechanische Verschlusssysteme spielen in der Konstruktionstechnik eine zentrale Rolle.
Klassische Verschlüsse basieren häufig auf einfachen Rast- oder Bajonettsystemen, die
jedoch in sicherheitskritischen Anwendungen unzureichend sein können, da sie
single-point-of-failure-Schwachstellen aufweisen.

Die vorliegende Arbeit beschreibt und analysiert einen \emph{doppelten Ringverschluss}
(engl. \textit{Double Ring Lock Mechanism}, kurz DRLM), der durch eine Zwei-Ebenen-Geometrie
mit jeweils neun Pins und neun Verschluss-Narben ein mechanisch redundantes und
hochsicheres Rastsystem realisiert. Der DRLM zeichnet sich durch folgende Merkmale aus:

\begin{itemize}[noitemsep]
  \item Zwei sequenzielle Rastebenen (obere Ebene $E_1$, untere Ebene $E_2$)
  \item Neun symmetrisch angeordnete Pins $\Pin = \{p_1, \ldots, p_9\}$
  \item Neun korrespondierende Verschluss-Narben $\Scar = \{s_1, \ldots, s_9\}$
  \item Rechtsdrehung zum ersten Einrasten, Linksdrehung zum zweiten Einrasten
  \item Entriegelung durch axialen Druckimpuls von oben
\end{itemize}

Das System wird vollständig geometrisch-formal beschrieben, physikalisch motiviert und
schließlich bezüglich seiner Optimalität nachgewiesen.

% ============================================================
\section{Systemarchitektur und Geometrie}
% ============================================================

\subsection{Grundlegende Definitionen}

\begin{definition}[Verschlusssystem DRLM]
  Ein \emph{Doppelter Ringverschluss} ist ein Tupel
  \[
    \text{DRLM} = (\Ring, \Pin, \Scar, E_1, E_2, \Phi_1, \Phi_2, \delta, \pi)
  \]
  wobei:
  \begin{itemize}[noitemsep]
    \item $\Ring$ der Verschlusskörper (Ringgeometrie) ist,
    \item $\Pin = \{p_1, \ldots, p_9\}$ die Menge der neun Pins bezeichnet,
    \item $\Scar = \{s_1, \ldots, s_9\}$ die Menge der neun Verschluss-Narben,
    \item $E_1, E_2 \subset \mathbb{R}^3$ die obere bzw. untere Rastebene,
    \item $\Phi_1 = +40°$ der Rechtsdrehwinkel zur ersten Einrastung,
    \item $\Phi_2 = -40°$ der Linksdrehwinkel zur zweiten Einrastung,
    \item $\delta > 0$ die axiale Drucktiefe für Zustand 2,
    \item $\pi > 0$ der Lösedruck (Entriegelungsimpuls).
  \end{itemize}
\end{definition}

\begin{definition}[Pinpositionen]
  Die neun Pins sind rotationssymmetrisch auf einem Kreis mit Radius $r_p$ angeordnet.
  Der Winkelabstand zweier benachbarter Pins beträgt
  \[
    \Delta\alpha = \frac{360°}{9} = 40°.
  \]
  Die Position des $k$-ten Pins in Polarkoordinaten ist:
  \[
    p_k = \left(r_p,\; \alpha_k\right), \quad \alpha_k = (k-1) \cdot 40°, \quad k = 1, \ldots, 9.
  \]
\end{definition}

\begin{definition}[Narben-Geometrie]
  Jede Verschluss-Narbe $s_k \in \Scar$ besitzt eine charakteristische L-förmige Geometrie
  auf der jeweiligen Rastebene:
  \[
    s_k = \left\{ \vvec{x} \in \mathbb{R}^2 \;\middle|\; \vvec{x} \in \Gamma_k^{(1)} \cup \Gamma_k^{(2)} \right\}
  \]
  mit dem vertikalen Einführkanal $\Gamma_k^{(1)}$ und dem horizontalen Rastkanal $\Gamma_k^{(2)}$.
\end{definition}

\subsection{Zwei-Ebenen-Modell}

Das DRLM operiert auf zwei axialen Ebenen:

\begin{align}
  E_1 &: \quad z = z_1 \quad \text{(obere Rastebene)}, \\
  E_2 &: \quad z = z_2 < z_1 \quad \text{(untere Rastebene)}.
\end{align}

Der axiale Abstand zwischen den Ebenen ist:
\[
  \Delta z = z_1 - z_2 > 0.
\]

\begin{figure}[H]
  \centering
  \begin{tikzpicture}[scale=1.0]
    % Ebene E1 (oben)
    \fill[blue!10, draw=blue!50, rounded corners=4pt]
      (-4.5, 3.5) rectangle (4.5, 5.5);
    \node[blue!70, font=\bfseries] at (0, 5.2) {Obere Rastebene $E_1$};

    % Ebene E2 (unten)
    \fill[green!10, draw=green!50, rounded corners=4pt]
      (-4.5, 0.0) rectangle (4.5, 2.0);
    \node[green!70, font=\bfseries] at (0, 1.7) {Untere Rastebene $E_2$};

    % Abstandspfeil
    \draw[<->, thick, gray] (4.8, 1.0) -- (4.8, 4.5);
    \node[right, gray] at (5.0, 2.75) {$\Delta z$};

    % Verschlusskörper
    \fill[ringcolor!40, draw=ringcolor, thick] (-0.5, 2.0) rectangle (0.5, 3.5);
    \node[ringcolor!80!black] at (0, 2.75) {$\Ring$};

    % Pins E1
    \foreach \x in {-3.5,-2.5,-1.5,-0.5,0.5,1.5,2.5,3.5,-4.0} {
      % Nur 9 Punkte anzeigen (vereinfacht)
    }
    \foreach \i in {1,...,9} {
      \pgfmathsetmacro{\xx}{-4.0 + (\i-1)*1.0}
      \pgfmathsetmacro{\yy}{4.5}
      \fill[pincolor] (\xx, \yy) circle (5pt);
      \node[above, pincolor, font=\tiny] at (\xx, \yy+0.1) {$p_{\i}$};
    }

    % Narben E2
    \foreach \i in {1,...,9} {
      \pgfmathsetmacro{\xx}{-4.0 + (\i-1)*1.0}
      \pgfmathsetmacro{\yy}{0.5}
      \draw[scarcolor, thick] (\xx-0.15, \yy+0.3) rectangle (\xx+0.15, \yy-0.1);
      \node[below, scarcolor, font=\tiny] at (\xx, \yy-0.15) {$s_{\i}$};
    }

    % Drehpfeile
    \draw[->, ultra thick, lockcolor!80!black, bend left=30]
      (1.5, 3.8) to node[above, font=\small]{Rechtsdrehung $\Phi_1$} (3.0, 3.8);
    \draw[->, ultra thick, presscolor!80!black, bend right=30]
      (3.0, 0.8) to node[below, font=\small]{Linksdrehung $\Phi_2$} (1.5, 0.8);

    % Druckpfeil
    \draw[->, ultra thick, red!70, line width=2pt] (0, 6.5) -- (0, 5.5);
    \node[red!70, font=\bfseries, above] at (0, 6.5) {Axialdruck $\pi$};

  \end{tikzpicture}
  \caption{Schematische Darstellung des Zwei-Ebenen-Modells des DRLM mit Pins ($p_i$, blau) und Narben ($s_i$, rot).}
  \label{fig:zwei_ebenen}
\end{figure}

% ============================================================
\section{Verschlusssequenz: Formale Beschreibung}
% ============================================================

\subsection{Zustandsraum}

\begin{definition}[Systemzustände]
  Das DRLM befindet sich zu jedem Zeitpunkt in einem der folgenden Zustände:
  \[
    \mathcal{Z} = \{Z_0, Z_1, Z_2, Z_E\}
  \]
  mit:
  \begin{itemize}[noitemsep]
    \item $Z_0$: Ausgangszustand (offen, Verschluss aufgesetzt, nicht gedreht)
    \item $Z_1$: Erster Rastzustand (Einrastung auf $E_1$ nach Rechtsdrehung)
    \item $Z_2$: Zweiter Rastzustand — vollständig geschlossen (Einrastung auf $E_1 \cup E_2$)
    \item $Z_E$: Entriegelter Zustand (nach Druckimpuls, beide Ebenen freigegeben)
  \end{itemize}
\end{definition}

\begin{definition}[Ãœbergangsfunktionen]
  Die Zustandsübergänge sind definiert durch:
  \begin{align}
    \tau_1 &: Z_0 \xrightarrow{+\Phi_1} Z_1 && \text{(Rechtsdrehung, erste Einrastung)}\\
    \tau_2 &: Z_1 \xrightarrow{\delta \downarrow,\; -\Phi_2} Z_2 && \text{(Axialdruck + Linksdrehung, zweite Einrastung)}\\
    \tau_3 &: Z_2 \xrightarrow{\pi \downarrow} Z_E && \text{(Zentraldruckimpuls, Entriegelung)}\\
    \tau_4 &: Z_E \xrightarrow{} Z_0 && \text{(Verschluss abnehmbar)}
  \end{align}
\end{definition}

Das vollständige Zustandsdiagramm ist in Abbildung~\ref{fig:zustandsdiagramm} dargestellt.

\begin{sidewaysfigure}
  \centering
  \begin{tikzpicture}[
    node distance=3.5cm,
    state/.style={draw, rounded rectangle, thick, minimum width=2.5cm, minimum height=1cm,
                  fill=white, drop shadow},
    >=Stealth
  ]
    \node[state, fill=unlockcolor!60] (Z0) {$Z_0$ --- offen};
    \node[state, fill=lockcolor!60, right=of Z0] (Z1) {$Z_1$ --- 1.~Einrastung};
    \node[state, fill=scarcolor!40, right=of Z1] (Z2) {$Z_2$ --- vollst.\ gesperrt};
    \node[state, fill=presscolor!40, below=2cm of Z2] (ZE) {$Z_E$ --- entriegelt};

    \draw[->, thick, blue!70] (Z0) -- node[above, font=\small]{$+\Phi_1 = +40°$} (Z1);
    \draw[->, thick, green!60!black] (Z1) -- node[above, font=\small]{$\delta{\downarrow}$, $-\Phi_2 = -40°$} (Z2);
    \draw[->, thick, red!70] (Z2) -- node[right, font=\small]{$\pi{\downarrow}$ (Zentraldruck)} (ZE);
    \draw[->, thick, gray] (ZE) -| node[below, font=\small, pos=0.25]{Abnahme} (Z0);
  \end{tikzpicture}
  \caption{Zustandsdiagramm des DRLM. Pfeile beschriften die jeweilige Auslösebedingung.}
  \label{fig:zustandsdiagramm}
\end{sidewaysfigure}

\subsection{Phase 1: Aufsetzen und erste Einrastung ($Z_0 \to Z_1$)}

\begin{figure}[H]
  \centering
  \begin{tikzpicture}[scale=1.2]
    % Grundplatte
    \fill[gray!20, draw=gray!60, rounded corners=3pt] (-3,-3) rectangle (3,3);
    \node[gray!60, font=\small] at (0, -2.7) {Gegenstück (Sockel)};

    % Neun Narben auf E1 (Einführöffnungen, radial verteilt)
    \foreach \ang in {0,40,80,120,160,200,240,280,320} {
      \pgfmathsetmacro{\rx}{2.1*cos(\ang)}
      \pgfmathsetmacro{\ry}{2.1*sin(\ang)}
      % Einführkanal (vertikal, oben)
      \fill[scarcolor!80, draw=scarcolor, rounded corners=1pt]
        (\rx-0.12, \ry-0.35) rectangle (\rx+0.12, \ry+0.35);
      \pgfmathtruncatemacro{\idx}{int(\ang/40)+1}
      \node[font=\tiny, black] at (\rx, \ry) {$s_{\idx}$};
    }

    % Verschlusskörper (Querschnitt)
    \fill[ringcolor!30, draw=ringcolor, thick, rounded corners=5pt]
      (-0.6,-0.6) rectangle (0.6, 0.6);
    \node[ringcolor!80!black, font=\bfseries\small] at (0,0) {$\Ring$};

    % Neun Pins (am Verschlusskörper, ausgefahren)
    \foreach \ang in {0,40,80,120,160,200,240,280,320} {
      \pgfmathsetmacro{\rx}{2.1*cos(\ang)}
      \pgfmathsetmacro{\ry}{2.1*sin(\ang)}
      \fill[pincolor, draw=pincolor!80!black] (\rx,\ry) circle (5pt);
      \pgfmathtruncatemacro{\idx}{int(\ang/40)+1}
    }

    % Drehpfeil
    \draw[->, ultra thick, lockcolor!80!black, bend left]
      (1.5, -1.5) arc[start angle=-45, end angle=45, radius=2.1];
    \node[lockcolor!80!black, font=\bfseries] at (2.8, 0) {$+40°$};

    \node[font=\bfseries, above] at (0, 3.2) {Phase 1: Aufsetzen + Rechtsdrehung ($Z_0 \to Z_1$)};

    % Druckpfeil (axial, von oben)
    \draw[->, thick, gray!80] (0, 1.5) -- (0, 0.65);
    \node[gray!80, font=\small, right] at (0.1, 1.1) {axial aufsetzen};
  \end{tikzpicture}
  \caption{Draufsicht auf Ebene $E_1$: Die neun Pins (blau) befinden sich über den Einführkanälen der Narben. Nach dem Aufsetzen erfolgt die Rechtsdrehung um $\Phi_1 = 40°$.}
  \label{fig:phase1}
\end{figure}

\subsubsection{Mechanisches Modell der ersten Einrastung}

Beim Aufsetzen greifen alle neun Pins $p_k$ in die vertikalen Einführkanäle $\Gamma_k^{(1)}$ der
korrespondierenden Narben $s_k$ ein. Durch die Rechtsdrehung um $\Phi_1 = 40°$ gleiten die Pins
in die horizontalen Rastkanäle $\Gamma_k^{(2)}$:

\begin{equation}
  p_k^{(1)} = \left(r_p,\; \alpha_k + \Phi_1\right) \quad \forall k \in \{1,\ldots,9\},
\end{equation}

wobei die Rastbedingung auf $E_1$ erfüllt ist, sobald:
\begin{equation}
  \label{eq:rastbedingung1}
  \forall k: \quad p_k^{(1)} \in s_k \cap \Gamma_k^{(2)}.
\end{equation}

Die Federkraft der Rastnasen bewirkt eine Haltekraft:
\begin{equation}
  F_{\text{Rast},1} = \sum_{k=1}^{9} f_k^{(1)} \cdot \cos(\beta_k),
\end{equation}
wobei $f_k^{(1)}$ die individuelle Federkraft der $k$-ten Rastnase und $\beta_k$ der
Eingriffswinkel ist. Für symmetrische Konfiguration gilt $\beta_k = \beta$ und $f_k^{(1)} = f$, sodass:
\begin{equation}
  F_{\text{Rast},1} = 9 \cdot f \cdot \cos(\beta).
\end{equation}

\subsection{Phase 2: Zweite Einrastung ($Z_1 \to Z_2$)}

\begin{figure}[H]
  \centering
  \begin{subfigure}[b]{0.48\textwidth}
    \centering
    \begin{tikzpicture}[scale=1.0]
      % Seitenansicht Verschluss
      % Ebene E1
      \fill[blue!15, draw=blue!50, rounded corners=3pt] (-2.8, 2.5) rectangle (2.8, 3.5);
      \node[blue!70, font=\small\bfseries] at (0, 3.2) {$E_1$ (1. Einrastung aktiv)};

      % Ebene E2
      \fill[green!15, draw=green!50, rounded corners=3pt] (-2.8, 0.0) rectangle (2.8, 1.0);
      \node[green!70, font=\small\bfseries] at (0, 0.7) {$E_2$ (bereit)};

      % Verschlusskörper
      \fill[ringcolor!30, draw=ringcolor, thick] (-0.5, 1.0) rectangle (0.5, 2.5);
      \node[ringcolor!80!black] at (0, 1.75) {$\Ring$};

      % Pins auf E1 (eingerastet)
      \foreach \x in {-2.2, -1.1, 0.0, 1.1, 2.2} {
        \fill[pincolor!70] (\x, 3.0) circle (4pt);
        \draw[pincolor!70, thick] (\x, 2.5) -- (\x, 3.0);
      }

      % Axialdruck-Pfeil
      \draw[->, ultra thick, red!70, line width=2pt] (0, 4.5) -- (0, 3.6);
      \node[red!70, right, font=\bfseries\small] at (0.1, 4.1) {$\delta$};

      % Drehpfeil
      \draw[->, ultra thick, presscolor!80!black] (1.8, 1.5) arc[start angle=0, end angle=120, radius=0.5];
      \node[presscolor!80!black, font=\bfseries\small, right] at (2.0, 1.7) {$-40°$};

      \node[font=\bfseries, above] at (0, 4.8) {Phase 2: Druck + Links\-dreh\-ung};
    \end{tikzpicture}
    \caption{Seitenansicht: axialer Druckimpuls $\delta$ und Linksdrehung $\Phi_2 = -40°$.}
  \end{subfigure}
  \hfill
  \begin{subfigure}[b]{0.48\textwidth}
    \centering
    \begin{tikzpicture}[scale=1.0]
      % Draufsicht E2 nach zweiter Einrastung
      \fill[gray!20, draw=gray!60, rounded corners=3pt] (-3,-3) rectangle (3,3);
      \node[gray!60, font=\small] at (0, -2.7) {Ebene $E_2$};

      % Narben E2 (L-förmig, gespiegelt zu E1)
      \foreach \ang in {0,40,80,120,160,200,240,280,320} {
        \pgfmathsetmacro{\rx}{2.1*cos(\ang)}
        \pgfmathsetmacro{\ry}{2.1*sin(\ang)}
        \pgfmathsetmacro{\rxn}{2.1*cos(\ang-20)}
        \pgfmathsetmacro{\ryn}{2.1*sin(\ang-20)}
        % L-Form: horizontal Rastkanal
        \fill[scarcolor!60, draw=scarcolor, rounded corners=1pt]
          (\rx-0.35, \ry-0.12) rectangle (\rx+0.15, \ry+0.12);
        \fill[scarcolor!60, draw=scarcolor, rounded corners=1pt]
          (\rx-0.12, \ry-0.35) rectangle (\rx+0.12, \ry+0.12);
        \pgfmathtruncatemacro{\idx}{int(\ang/40)+1}
        \node[font=\tiny, black] at (\rx+0.3, \ry+0.3) {$s_{\idx}'$};
      }

      % Pins nach zweiter Einrastung
      \foreach \ang in {0,40,80,120,160,200,240,280,320} {
        \pgfmathsetmacro{\rx}{2.1*cos(\ang)}
        \pgfmathsetmacro{\ry}{2.1*sin(\ang)}
        \fill[pincolor, draw=pincolor!80!black] (\rx,\ry) circle (5pt);
      }

      % Einraststatus
      \foreach \ang in {0,40,80,120,160,200,240,280,320} {
        \pgfmathsetmacro{\rx}{2.1*cos(\ang)*0.55}
        \pgfmathsetmacro{\ry}{2.1*sin(\ang)*0.55}
        \node[green!60!black, font=\bfseries\Large] at (\rx,\ry) {};
      }

      % Zentrum
      \fill[presscolor!50, draw=presscolor, thick] (0,0) circle (0.5);
      \node[presscolor!80!black, font=\bfseries\small] at (0,0) {$\pi$};
      \draw[->, ultra thick, red!70] (0, 1.8) -- (0, 0.55);
      \node[red!70, font=\bfseries\small, right] at (0.1, 1.2) {Lösedruck};

      \node[font=\bfseries, above] at (0, 3.2) {$Z_2$: Vollständig eingerastet};
    \end{tikzpicture}
    \caption{Draufsicht auf $E_2$ nach vollständiger zweiter Einrastung ($Z_2$). Zentrum: Entriegelungspunkt.}
  \end{subfigure}
  \caption{Phase 2: Zweite Einrastung durch Axialdruck und Linksdrehung. Links: Seitenansicht der Betätigung. Rechts: Endposition auf $E_2$.}
  \label{fig:phase2}
\end{figure}

\subsubsection{Mechanisches Modell der zweiten Einrastung}

In Zustand $Z_1$ wird der Verschlusskörper axial um $\delta$ abgesenkt, sodass die
Pins in die Ebene $E_2$ eintreten. Gleichzeitig erfolgt eine Linksdrehung um $|\Phi_2| = 40°$:

\begin{align}
  z: \quad & z_1 \to z_1 - \delta = z_2 \\
  \alpha: \quad & \alpha_k^{(1)} \to \alpha_k^{(1)} - |\Phi_2| = \alpha_k^{(2)}
\end{align}

Die kombinierte Gesamtdrehung der Pins bezüglich der Ausgangsstellung beträgt:
\begin{equation}
  \Delta\alpha_{\text{ges}} = \Phi_1 + (-|\Phi_2|) = 40° - 40° = 0°.
\end{equation}

Dies bedeutet, dass die Pins auf $E_2$ in ihrer ursprünglichen Winkelposition einrasten —
ein wesentliches Merkmal für die Symmetrie und Reproduzierbarkeit des Systems.

Die Gesamthaltekraft im vollständig gesperrten Zustand $Z_2$ ist:
\begin{equation}
  \boxed{F_{\text{Rast,ges}} = F_{\text{Rast},1}^{(E_1)} + F_{\text{Rast},2}^{(E_2)} = 9 f^{(1)}\cos\beta + 9 f^{(2)}\cos\beta = 9\cos\beta\left(f^{(1)} + f^{(2)}\right).}
\end{equation}

\subsection{Entriegelung durch Zentraldruckimpuls ($Z_2 \to Z_E$)}

\begin{figure}[H]
  \centering
  \begin{tikzpicture}[scale=1.1]
    % Hintergrund
    \fill[gray!10, draw=gray!40, rounded corners=5pt] (-4.5,-1) rectangle (4.5, 5.5);

    % Ebene E1
    \fill[blue!15, draw=blue!50, rounded corners=3pt] (-4.0, 3.5) rectangle (4.0, 4.5);
    \node[blue!70, font=\bfseries\small] at (0, 4.2) {$E_1$ — wird freigegeben};

    % Ebene E2
    \fill[green!15, draw=green!50, rounded corners=3pt] (-4.0, 1.0) rectangle (4.0, 2.0);
    \node[green!70, font=\bfseries\small] at (0, 1.7) {$E_2$ — wird freigegeben};

    % Verschlusskörper
    \fill[ringcolor!30, draw=ringcolor, thick] (-0.6, 2.0) rectangle (0.6, 3.5);
    \node[ringcolor!80!black, font=\bfseries] at (0, 2.75) {$\Ring$};

    % Zentraler Druckknopf
    \fill[presscolor!60, draw=presscolor, very thick] (-0.4, 3.5) rectangle (0.4, 4.2);
    \node[presscolor!80!black, font=\bfseries\tiny] at (0, 3.85) {Zentrum};

    % Druckpfeil
    \draw[->, ultra thick, red!80, line width=2.5pt] (0, 5.5) -- (0, 4.2);
    \node[red!80, font=\bfseries, right] at (0.15, 5.0) {$F_\pi$ (Lösedruck)};

    % Freigabepfeile E1 (Pins bewegen sich raus)
    \foreach \x in {-3.2, -2.1, -1.0, 1.0, 2.1, 3.2} {
      \draw[->, thick, blue!60] (\x, 4.0) -- (\x+0.0, 4.8);
    }

    % Freigabepfeile E2
    \foreach \x in {-3.2, -2.1, -1.0, 1.0, 2.1, 3.2} {
      \draw[->, thick, green!60!black] (\x, 1.5) -- (\x+0.0, 0.5);
    }

    % Beschriftung Freigabe
    \node[blue!60, font=\small, left] at (-4.0, 4.4) {Pins rasten aus $E_1$};
    \node[green!60!black, font=\small, left] at (-4.0, 0.3) {Pins rasten aus $E_2$};

    % Gesamt-Entriegelung
    \node[red!70, font=\bfseries, below] at (0, -0.7)
      {Vollständige Entriegelung: $Z_2 \to Z_E$ bei $F_\pi > F_{\text{Rast,ges}}$};
  \end{tikzpicture}
  \caption{Entriegelungsphase: Ein einzelner axialer Druckimpuls $F_\pi$ auf das Zentrum des Verschlusses löst simultan beide Rastebenen $E_1$ und $E_2$.}
  \label{fig:entriegelung}
\end{figure}

\subsubsection{Kinematik der Entriegelung}

Der Zentraldruckknopf ist mit einer Koppelmechanik verbunden, die beide Rastebenen
simultan ansteuert. Die Entriegelungskraft muss die Gesamthaltekraft überwinden:

\begin{equation}
  F_\pi > F_{\text{Rast,ges}} = 9\cos\beta\left(f^{(1)} + f^{(2)}\right).
\end{equation}

Durch die axiale Ãœbertragung wirkt $F_\pi$ gleichzeitig auf beide Ebenen, was die
charakteristische Einknopf-Bedienung ermöglicht.

% ============================================================
\section{Formaler Nachweis der Optimalität}
% ============================================================

\subsection{Optimalitätskriterien}

Ein mechanischer Verschluss wird als \emph{optimal} bezeichnet, wenn er folgende Kriterien
gleichzeitig erfüllt:

\begin{enumerate}[label=\textbf{K\arabic*}]
  \item \textbf{Maximale Rastredundanz}: Jeder Einrastzustand wird durch mehrere unabhängige
        Eingriffspunkte gesichert.
  \item \textbf{Minimale Fehlbedienbarkeit}: Unbeabsichtigtes Öffnen ist durch mehr als eine
        unabhängige Aktion ausgeschlossen.
  \item \textbf{Symmetrische Kraftverteilung}: Die Haltekräfte verteilen sich gleichmäßig über
        den gesamten Umfang.
  \item \textbf{Eindeutige Entriegelung}: Exakt eine definierte Aktion führt zur Freigabe.
  \item \textbf{Sequenzielle Sicherheit}: Das vollständige Schließen erfordert zwei unabhängige,
        geordnete Aktionen.
\end{enumerate}

\subsection{Nachweis der 9-Pin-Optimalität}

\begin{theorem}[Optimale Pinanzahl]
  \label{thm:pinoptimal}
  Für einen kreisförmigen Doppelring-Verschluss mit Drehwinkeln von $\Phi_1 = 40°$ und
  $|\Phi_2| = 40°$ ist die Anzahl $n = 9$ Pins die eindeutig optimale Konfiguration.
\end{theorem}

\begin{proof}
  Die Pinanzahl $n$ muss folgende Bedingungen gleichzeitig erfüllen:

  \textbf{(i) Ganzzahliger Teiler der Kreisteilung:}
  Der Winkelabstand der Pins muss mit dem Drehwinkel harmonieren:
  \[
    \Delta\alpha = \frac{360°}{n} \overset{!}{=} \Phi_1 = |\Phi_2|.
  \]
  Für $\Phi_1 = 40°$ folgt:
  \[
    n = \frac{360°}{40°} = 9. \quad \checkmark
  \]

  \textbf{(ii) Vollständige Gleichverteilung:}
  Nur für $n = 9$ sind alle Pins und alle Narben bei Drehung um $\Phi_1$ gleichzeitig in
  Eingriffsposition, da $9 \cdot 40° = 360°$ eine vollständige Kreisteilung ergibt.

  \textbf{(iii) Rückkehreigenschaft:}
  Die Gesamtdrehung $\Phi_1 - |\Phi_2| = 0°$ führt die Pins auf $E_2$ in die
  Ursprungsposition zurück. Dies ist nur möglich, wenn $\Phi_1 = |\Phi_2|$ und
  $n = 360°/\Phi_1 \in \mathbb{N}$. Für $n=9$ gilt:
  \[
    n \cdot \Phi_1 = 9 \cdot 40° = 360°, \quad n \cdot |\Phi_2| = 9 \cdot 40° = 360°. \quad \checkmark
  \]

  Jede andere Pinanzahl $n' \neq 9$ mit $\Phi_1 = 40°$ würde eine nicht-ganzzahlige
  Kreisteilung oder asymmetrische Kraftverteilung erzeugen. Daher ist $n = 9$ eindeutig optimal.
\end{proof}

\begin{theorem}[Redundante Sicherheit durch Zwei-Ebenen-System]
  \label{thm:redundanz}
  Das DRLM erfüllt Kriterium \textbf{K2} (Minimale Fehlbedienbarkeit): Für eine
  unbeabsichtigte Öffnung wären mindestens zwei unabhängige, sequenzielle Fehlaktionen
  erforderlich.
\end{theorem}

\begin{proof}
  Im Zustand $Z_2$ ist der Verschluss auf beiden Ebenen $E_1$ und $E_2$ aktiv. Die
  Überführung $Z_2 \to Z_0$ erfordert:
  \begin{enumerate}
    \item Aufbringen eines Druckimpulses $F_\pi > F_{\text{Rast,ges}}$ im Zentrum
          (zur Überführung $Z_2 \to Z_E$), und
    \item Axiales Abheben des Verschlusskörpers (zur Überführung $Z_E \to Z_0$).
  \end{enumerate}
  Diese beiden Aktionen sind physikalisch unabhängig und können nicht durch eine einzige
  zufällige Kraft ausgelöst werden, da $F_\pi$ ausschließlich als zentrale Axialkraft wirkt,
  während die Abnahme eine andere Kraft\-richtung erfordert. Damit folgt K2. $\square$
\end{proof}

\begin{theorem}[Symmetrische Kraftverteilung]
  \label{thm:symmetrie}
  Im Zustand $Z_2$ ist die Haltekraft $F_{\text{Rast,ges}}$ gleichmäßig auf alle $n = 9$
  Pins und Narben verteilt. Das Drehmoment um die Systemachse ist null.
\end{theorem}

\begin{proof}
  Die Pins sind bei Winkelabstand $\Delta\alpha = 40°$ gleichmäßig auf dem Kreis mit
  Radius $r_p$ angeordnet. Die von jedem Pin aufgebrachte Radialkraft ist $f \cos\beta$ (symmetrisch).
  Das resultierende Drehmoment ergibt sich als Vektorsumme:
  \[
    M = \sum_{k=1}^{9} r_p \cdot f\cos\beta \cdot \hat{e}_{\perp,k}
      = r_p f\cos\beta \sum_{k=1}^{9} \hat{e}_{\perp,k}.
  \]
  Da die Einheitsvektoren $\hat{e}_{\perp,k}$ eine vollständige Kreisteilung bilden, gilt:
  \[
    \sum_{k=1}^{9} \hat{e}_{\perp,k} = \vvec{0}.
  \]
  Damit ist $M = 0$, d.\,h. der Verschluss ist drehmomentfrei gehalten. $\square$
\end{proof}

\begin{corollary}[Vibrationsresistenz]
  Die drehmomentfreie Halterung (Satz~\ref{thm:symmetrie}) impliziert, dass externe
  Vibrationen kein Nettodrehmoment auf den Verschluss ausüben können und somit keine
  Selbstentriegelung induzieren.
\end{corollary}

\subsection{Nachweis der Eindeutigkeit des Lösemechanismus}

\begin{theorem}[Eindeutige Entriegelung]
  \label{thm:einheit}
  Im Zustand $Z_2$ existiert genau eine physikalische Aktion, die zur Entriegelung führt:
  der axiale Zentraldruckimpuls $F_\pi > F_{\text{Rast,ges}}$.
\end{theorem}

\begin{proof}
  Wir betrachten alle möglichen Kraftkomponenten auf den Verschlusskörper $\Ring$:
  \begin{itemize}
    \item \textit{Axiale Zugkraft} (Abheben): blockiert durch die Rastnasen beider Ebenen;
          Entriegelung erfordert $F_{\text{zug}} > F_{\text{Rast,ges}}$ — jedoch löst diese
          Kraft nicht den Druckmechanismus aus, da sie in umgekehrter Richtung wirkt.
    \item \textit{Drehmoment}: führt zu höherer Klemmkraft durch Geometrie der L-Narben;
          Drehung kann den Zustand $Z_2$ nicht verlassen.
    \item \textit{Laterale Kraft}: durch Kreisgeometrie in radialer Richtung blockiert.
    \item \textit{Axialer Zentraldruck} ($F_\pi$): aktiviert die Koppelmechanik und löst
          gleichzeitig beide Ebenen.
  \end{itemize}
  Da keine andere Kraftkomponente die geometrische Verriegelung überwinden kann, ist
  der Zentraldruckimpuls die einzige Entriegelungsaktion. $\square$
\end{proof}

% ============================================================
\section{Gesamtdarstellung: TikZ-Ãœbersicht}
% ============================================================

\begin{figure}[H]
  \centering
  \begin{tikzpicture}[scale=0.9]
    %--- Z0: Aufsetzen ---
    \begin{scope}[xshift=-7cm, yshift=0cm]
      \fill[gray!15, draw=gray!50, rounded corners=3pt] (-1.5,-1.5) rectangle (1.5,1.5);
      \foreach \ang in {0,40,80,120,160,200,240,280,320} {
        \pgfmathsetmacro{\rx}{1.1*cos(\ang)}
        \pgfmathsetmacro{\ry}{1.1*sin(\ang)}
        \fill[unlockcolor!80, draw=gray] (\rx,\ry) circle (3.5pt);
      }
      \fill[ringcolor!30, draw=ringcolor, thick] (-0.3,-0.3) rectangle (0.3,0.3);
      \node[font=\bfseries\small, above] at (0,1.7) {$Z_0$: Aufgesetzt};
      \draw[->, thick, blue!70] (0,-2.0) -- (0,-1.6);
      \node[blue!70, font=\tiny, below] at (0,-2.1) {axial absenken};
    \end{scope}

    %--- Z1: Erste Einrastung ---
    \begin{scope}[xshift=-2.5cm, yshift=0cm]
      \fill[gray!15, draw=gray!50, rounded corners=3pt] (-1.5,-1.5) rectangle (1.5,1.5);
      \foreach \ang in {0,40,80,120,160,200,240,280,320} {
        \pgfmathsetmacro{\rx}{1.1*cos(\ang)}
        \pgfmathsetmacro{\ry}{1.1*sin(\ang)}
        \fill[pincolor, draw=pincolor!80!black] (\rx,\ry) circle (3.5pt);
        % Rastmarkierung
        \node[lockcolor!80!black, font=\tiny] at (\rx*0.6, \ry*0.6) {};
      }
      \fill[ringcolor!30, draw=ringcolor, thick] (-0.3,-0.3) rectangle (0.3,0.3);
      % Rastung E1
      \draw[lockcolor!80!black, very thick] (0,0) circle (1.15);
      \node[font=\bfseries\small, above] at (0,1.7) {$Z_1$: 1. Einrastung};
      \node[lockcolor!80!black, font=\small\bfseries] at (0,-1.85) {$E_1$ aktiv};
      % Drehpfeil
      \draw[->, thick, lockcolor!80!black] (1.3,0) arc[start angle=0,end angle=40,radius=1.3];
      \node[lockcolor!80!black, font=\tiny, right] at (1.7, 0.4) {$+40°$};
    \end{scope}

    %--- Z2: Vollständig gesperrt ---
    \begin{scope}[xshift=2.5cm, yshift=0cm]
      \fill[gray!15, draw=gray!50, rounded corners=3pt] (-1.5,-1.5) rectangle (1.5,1.5);
      \foreach \ang in {0,40,80,120,160,200,240,280,320} {
        \pgfmathsetmacro{\rx}{1.1*cos(\ang)}
        \pgfmathsetmacro{\ry}{1.1*sin(\ang)}
        \fill[pincolor, draw=pincolor!80!black] (\rx,\ry) circle (3.5pt);
      }
      \fill[ringcolor!50, draw=ringcolor, thick] (-0.3,-0.3) rectangle (0.3,0.3);
      % Doppelrastung
      \draw[scarcolor, very thick, double] (0,0) circle (1.15);
      % Zentralknopf
      \fill[presscolor!60, draw=presscolor] (0,0) circle (0.28);
      \node[font=\bfseries\small, above] at (0,1.7) {$Z_2$: Vollg. gesperrt};
      \node[scarcolor, font=\small\bfseries] at (0,-1.85) {$E_1 + E_2$ aktiv};
      % Drehpfeil
      \draw[->, thick, presscolor!80!black] (1.3, 0.3) arc[start angle=20,end angle=-25,radius=1.3];
      \node[presscolor!80!black, font=\tiny, right] at (1.7, 0) {$-40°$};
    \end{scope}

    %--- ZE: Entriegelt ---
    \begin{scope}[xshift=7cm, yshift=0cm]
      \fill[gray!15, draw=gray!50, rounded corners=3pt] (-1.5,-1.5) rectangle (1.5,1.5);
      \foreach \ang in {0,40,80,120,160,200,240,280,320} {
        \pgfmathsetmacro{\rx}{1.1*cos(\ang)}
        \pgfmathsetmacro{\ry}{1.1*sin(\ang)}
        \fill[unlockcolor!80, draw=gray] (\rx,\ry) circle (3.5pt);
      }
      \fill[ringcolor!20, draw=gray, thick] (-0.3,-0.3) rectangle (0.3,0.3);
      \fill[presscolor!30, draw=presscolor!60] (0,0) circle (0.28);
      \draw[->, ultra thick, red!70] (0, 0.28) -- (0, 1.1);
      \node[font=\bfseries\small, above] at (0,1.7) {$Z_E$: Entriegelt};
      \node[red!70, font=\small\bfseries] at (0,-1.85) {Druck $F_\pi$};
    \end{scope}

    %--- Ãœbergangspfeile ---
    \draw[->, ultra thick, blue!70] (-5.4,0) -- (-4.1,0)
      node[midway, above, font=\small\bfseries] {$+40°$};
    \draw[->, ultra thick, green!60!black] (-0.9,0) -- (0.9,0)
      node[midway, above, font=\small\bfseries] {$\delta{\downarrow}, -40°$};
    \draw[->, ultra thick, red!70] (4.1,0) -- (5.4,0)
      node[midway, above, font=\small\bfseries] {$F_\pi$};

    % Legende
    \begin{scope}[xshift=-7cm, yshift=-3cm]
      \fill[pincolor] (0,0) circle (4pt); \node[right, font=\small] at (0.2,0) {Pin (eingerastet)};
      \fill[unlockcolor!80, draw=gray] (2.5,0) circle (4pt); \node[right, font=\small] at (2.7,0) {Pin (frei)};
      \fill[presscolor!60, draw=presscolor] (5.0,0) circle (4pt); \node[right, font=\small] at (5.2,0) {Entriegelungszentrum};
    \end{scope}
  \end{tikzpicture}
  \caption{Vollständige Zustandssequenz des DRLM. Von links nach rechts: Ausgangszustand $Z_0$,
           erste Einrastung $Z_1$ (Rechtsdrehung $+40°$), vollständige Verriegelung $Z_2$
           (Axialdruck + Linksdrehung $-40°$), Entriegelung $Z_E$ (Zentraldruckimpuls $F_\pi$).}
  \label{fig:gesamt}
\end{figure}

% ============================================================
\section{Parameterstudie und Kenngrößen}
% ============================================================

\subsection{Geometrische Parameter}

\begin{table}[H]
  \centering
  \caption{Charakteristische geometrische Kenngrößen des DRLM.}
  \label{tab:parameter}
  \begin{tabular}{llll}
    \toprule
    \textbf{Parameter} & \textbf{Symbol} & \textbf{Wert} & \textbf{Einheit} \\
    \midrule
    Anzahl Pins / Narben & $n$ & 9 & --- \\
    Winkelabstand (Pins) & $\Delta\alpha$ & 40 & \(^\circ\) \\
    Rechtsdrehwinkel & $\Phi_1$ & $+40$ & \(^\circ\) \\
    Linksdrehwinkel & $\Phi_2$ & $-40$ & \(^\circ\) \\
    Gesamtdrehung (auf $E_2$) & $\Phi_{\text{ges}}$ & 0 & \(^\circ\) \\
    Anzahl Rastebenen & $N_E$ & 2 & --- \\
    Axiale Drucktiefe & $\delta$ & $>0$ & mm \\
    Kreisteilung (gesamt) & $n \cdot \Delta\alpha$ & 360 & \(^\circ\) \\
    \bottomrule
  \end{tabular}
\end{table}

\subsection{Kraftanalyse}

Die Mindest-Entriegelungskraft skaliert linear mit der Federkonstante und der Pinanzahl:

\begin{equation}
  F_{\pi,\min} = 9\cos\beta \left(f^{(1)} + f^{(2)}\right) + F_{\text{Sicherheit}},
\end{equation}

wobei $F_{\text{Sicherheit}} > 0$ ein konstruktiver Sicherheitszuschlag ist.

\subsection{Vergleich: Einfacher vs. doppelter Ringverschluss}

\begin{table}[H]
  \centering
  \caption{Vergleich einfacher vs. doppelter Ringverschluss.}
  \label{tab:vergleich}
  \begin{tabular}{lcc}
    \toprule
    \textbf{Eigenschaft} & \textbf{Einfacher Ring} & \textbf{DRLM (doppelt)} \\
    \midrule
    Rastebenen & 1 & 2 \\
    Mindestverschlussschritte & 1 & 2 \\
    Haltekraft (relativ) & $9f\cos\beta$ & $9\cos\beta(f^{(1)}+f^{(2)})$ \\
    Entriegelungsschritte & 1 & 1 (zentral) \\
    Redundanz & gering & hoch \\
    Vibrationsresistenz & mittel & hoch (drehmomentfrei) \\
    Fehlbedienbarkeit & mittel & minimal \\
    \bottomrule
  \end{tabular}
\end{table}

% ============================================================
\section{Zusammenfassung und Schlussfolgerung}
% ============================================================

In dieser Arbeit wurde der \emph{Doppelte Ringverschluss (DRLM)} mit 9-Pin-Konfiguration
vollständig beschrieben, formal modelliert und hinsichtlich seiner Optimalität nachgewiesen.
Die wesentlichen Ergebnisse sind:

\begin{enumerate}
  \item \textbf{Geometrische Optimalität}: Die Anzahl von $n = 9$ Pins ergibt sich zwingend
        aus dem Drehwinkel $\Phi = 40°$ und der Kreissymmetrie (Satz~\ref{thm:pinoptimal}).

  \item \textbf{Zwei-Ebenen-Redundanz}: Durch die sequenzielle Verschlusssequenz
        $Z_0 \to Z_1 \to Z_2$ sind zwei unabhängige mechanische Aktionen zum vollständigen
        Schließen notwendig, was Fehlbedienung ausschließt (Satz~\ref{thm:redundanz}).

  \item \textbf{Drehmomentfreiheit}: Die symmetrische 9-fache Anordnung garantiert
        verschwindende Nettodrehmomente und damit Vibrationsresistenz (Satz~\ref{thm:symmetrie}).

  \item \textbf{Eindeutige Entriegelung}: Genau eine physikalische Aktion — der axiale
        Zentraldruckimpuls — führt zur simultanen Freigabe beider Ebenen (Satz~\ref{thm:einheit}).
\end{enumerate}

Der DRLM stellt damit eine mechanisch optimale Lösung für hochsichere, redundante
Verschlusssysteme dar, die besonders in sicherheitskritischen Anwendungen (Luft- und
Raumfahrt, Medizintechnik, Automotive) Vorteile gegenüber einfachen Rastsystemen bietet.

\newpage
\appendix

\section{Ergänzende TikZ-Illustration: Pin-Narben-Eingriff}

\begin{figure}[H]
  \centering
  \begin{tikzpicture}[scale=1.4]
    % L-förmige Narbe
    \fill[scarcolor!20, draw=scarcolor, thick, rounded corners=2pt]
      (-0.2, 0) -- (-0.2, 1.2) -- (0.2, 1.2) -- (0.2, 0.4)
      -- (1.0, 0.4) -- (1.0, 0) -- cycle;

    \node[scarcolor!80!black, font=\bfseries] at (0.4, -0.25) {Narbe $s_k$ (L-Form)};

    % Einführkanal
    \draw[<->, gray, dashed] (-0.22, 0.4) -- (-0.22, 1.2);
    \node[gray, font=\tiny, left] at (-0.22, 0.8) {$\Gamma_k^{(1)}$};

    % Rastkanal
    \draw[<->, gray, dashed] (0.2, 0.02) -- (1.0, 0.02);
    \node[gray, font=\tiny, below] at (0.6, 0.02) {$\Gamma_k^{(2)}$};

    % Pin in Einführposition
    \begin{scope}[xshift=3.0cm]
      \fill[scarcolor!20, draw=scarcolor, thick, rounded corners=2pt]
        (-0.2, 0) -- (-0.2, 1.2) -- (0.2, 1.2) -- (0.2, 0.4)
        -- (1.0, 0.4) -- (1.0, 0) -- cycle;
      \fill[pincolor, draw=pincolor!80!black] (0.0, 0.85) circle (0.17);
      \draw[->, thick, blue!70] (0.0, 1.2) -- (0.0, 0.57);
      \node[blue!70, font=\bfseries\small, right] at (0.2, 1.0) {$p_k$ absenken};
      \node[font=\bfseries\small, below] at (0.4, -0.25) {Aufsetzen ($Z_0$)};
    \end{scope}

    % Pin nach Rechtsdrehung (in Rastkanal)
    \begin{scope}[xshift=6.5cm]
      \fill[scarcolor!20, draw=scarcolor, thick, rounded corners=2pt]
        (-0.2, 0) -- (-0.2, 1.2) -- (0.2, 1.2) -- (0.2, 0.4)
        -- (1.0, 0.4) -- (1.0, 0) -- cycle;
      \fill[pincolor, draw=pincolor!80!black] (0.65, 0.2) circle (0.17);
      \draw[->, thick, lockcolor!80!black] (0.0, 0.2) -- (0.47, 0.2);
      \node[lockcolor!80!black, font=\bfseries\small, above] at (0.3, 0.37) {$\Phi_1 = +40°$};
      \fill[lockcolor!60, draw=lockcolor] (0.65, 0.2) circle (0.17);
      \node[font=\bfseries\small, below] at (0.4, -0.25) {Eingerastet ($Z_1$)};
    \end{scope}
  \end{tikzpicture}
  \caption{Detailansicht des Pin-Narben-Eingriffs: links die L-förmige Narbengeometrie,
           Mitte der Pin beim Absenken ($Z_0$), rechts der Pin im Rastkanal nach Rechtsdrehung ($Z_1$).}
  \label{fig:pinnarbe}
\end{figure}

\begin{figure}[H]
  \centering
  \begin{tikzpicture}[scale=1.3]
    % Vollständiger Kreis der 9 Pins mit Winkelbeschriftung
    \fill[gray!10, draw=gray!40] (0,0) circle (2.5cm);
    \draw[gray!50, dashed] (0,0) circle (2.1cm);
    \node[gray!70, font=\small] at (0,0) {Achse};

    \foreach \i in {1,...,9} {
      \pgfmathsetmacro{\ang}{(\i-1)*40}
      \pgfmathsetmacro{\rx}{2.1*cos(\ang)}
      \pgfmathsetmacro{\ry}{2.1*sin(\ang)}
      \fill[pincolor, draw=pincolor!80!black] (\rx,\ry) circle (5pt);
      \node[font=\tiny\bfseries, pincolor!80!black]
        at ({2.45*cos(\ang)},{2.45*sin(\ang)}) {$p_{\i}$};
      % Winkelmarkierung
      \draw[gray!50, thin] (0,0) -- (\rx,\ry);
      \node[font=\tiny, gray!70] at ({1.5*cos(\ang+20)},{1.5*sin(\ang+20)}) {$40°$};
    }

    % Drehpfeile
    \draw[->, ultra thick, lockcolor!80!black]
      ({2.3*cos(0)},{2.3*sin(0)}) arc[start angle=0, end angle=40, radius=2.3];
    \node[lockcolor!80!black, font=\bfseries\small] at (3.2, 0.5) {$\Phi_1 = +40°$};

    \draw[->, ultra thick, presscolor!80!black]
      ({2.3*cos(40)},{2.3*sin(40)}) arc[start angle=40, end angle=0, radius=2.3];
    \node[presscolor!80!black, font=\bfseries\small] at (3.2, -0.5) {$\Phi_2 = -40°$};

    \node[font=\bfseries, above] at (0, 2.75) {9-Pin-Kreisanordnung ($\Delta\alpha = 40°$)};
  \end{tikzpicture}
  \caption{Kreisanordnung der neun Pins mit Winkelabstand $\Delta\alpha = 40°$. Rechtsdrehung
           $\Phi_1 = +40°$ für erste Einrastung, Linksdrehung $\Phi_2 = -40°$ für zweite Einrastung.}
  \label{fig:kreis9pin}
\end{figure}


% ============================================================
\section{CAD-Konstruktionszeichnung und Fertigungsspezifikation}
% ============================================================

Dieses Kapitel enthält die vollständigen technischen Zeichnungen des Doppelten
Ringverschlusses (DRLM) in normgerechter Darstellung nach DIN~ISO~128 sowie alle für
die Fertigung und Montage notwendigen Maß-, Toleranz- und Oberflächenangaben.
Alle Maße sind in Millimeter angegeben, sofern nicht anders vermerkt.

% -----------------------------------------------------------
\subsection{Bauteilübersicht und Stückliste}
% -----------------------------------------------------------

Das DRLM besteht aus drei Hauptbaugruppen:

\begin{sidewaystable}
  \centering
  \caption{Stückliste DRLM -- Gesamtsystem}
  \label{tab:stueckliste}
  \renewcommand{\arraystretch}{1.3}
  \begin{tabular}{clllcc}
    \toprule
    \textbf{Pos.} & \textbf{Benennung} & \textbf{Werkstoff} & \textbf{Norm/Zeichnung} & \textbf{Menge} & \textbf{Masse [g]} \\
    \midrule
    1  & Verschlusskörper (Oberring)    & EN~1.4301 (V2A)       & DRG-001 & 1  & 48{,}2 \\
    2  & Sockelring (Unterring)         & EN~1.4301 (V2A)       & DRG-002 & 1  & 51{,}7 \\
    3  & Pin, zylindrisch $\varnothing$3{,}0    & EN~1.4571 (V4A)       & DRG-003 & 9  &  0{,}8 \\
    4  & Druckfeder (Pin-Feder)         & Federstahl 1.4310     & DIN~2098 & 9  &  0{,}1 \\
    5  & Verschluss-Narbenplatte        & EN~1.4301 (V2A)       & DRG-004 & 2  & 22{,}4 \\
    6  & Zentraldruckknopf              & EN~1.4301 (V2A)       & DRG-005 & 1  &  6{,}3 \\
    7  & Koppelstange (Entriegelung)    & EN~1.4571 (V4A)       & DRG-006 & 1  &  3{,}1 \\
    8  & Rückstellfeder (zentral)       & Federstahl 1.4310     & DIN~2098 & 1  &  0{,}4 \\
    9  & Sicherungsring $\varnothing$42        & Federstahl            & DIN~471  & 2  &  0{,}3 \\
    10 & O-Ring (Dichtung)              & FKM 75~Shore~A        & DIN~3771 & 2  &  0{,}2 \\
    \midrule
    \multicolumn{5}{l}{\textbf{Gesamtmasse (ohne Befestigungsmittel):}} & \textbf{$\approx$\,190 g} \\
    \bottomrule
  \end{tabular}
\end{sidewaystable}

% -----------------------------------------------------------
\subsection{Hauptabmessungen -- Gesamtansicht (Normschnitt)}
% -----------------------------------------------------------

\begin{figure}[H]
\centering
\begin{tikzpicture}[scale=1.0,
  cadline/.style={draw=black, line width=0.5pt},
  cadthick/.style={draw=black, line width=1.0pt},
  cadcenter/.style={draw=black!50, line width=0.3pt, dash pattern=on 4pt off 2pt on 1pt off 2pt},
  cadmass/.style={draw=black, line width=0.3pt},
  hatch/.style={pattern=north east lines, pattern color=black!40},
  dim/.style={<->, >=Stealth, thin, black},
  diml/.style={draw=black, thin},
]

%% ---- SCHNITTANSICHT: LÄNGSSCHNITT (vertikal, durch Mittelachse) ----

% Mittelachse
\draw[cadcenter] (0, -6.5) -- (0, 6.5);
\node[right, font=\tiny, black!50] at (0.05, 6.3) {Mittelachse};

%--- SOCKELRING (Pos. 2, unten) ---
% Außenkontur
\fill[hatch] (-3.0,-5.5) rectangle (-0.2,-2.5);
\fill[hatch] ( 0.2,-5.5) rectangle ( 3.0,-2.5);
\draw[cadthick] (-3.0,-5.5) rectangle (-0.2,-2.5);
\draw[cadthick] ( 0.2,-5.5) rectangle ( 3.0,-2.5);
% Boden
\fill[hatch] (-3.0,-6.0) rectangle (3.0,-5.5);
\draw[cadthick] (-3.0,-6.0) rectangle (3.0,-5.5);
% Bohrung innen oben
\draw[cadthick] (-0.2,-2.5) -- (-0.2,-1.8);
\draw[cadthick] ( 0.2,-2.5) -- ( 0.2,-1.8);
% Narbenplatte E2 (Pos. 5, untere Ebene)
\fill[gray!60] (-2.8,-3.3) rectangle (-0.25,-2.8);
\fill[gray!60] ( 0.25,-3.3) rectangle ( 2.8,-2.8);
\draw[cadline] (-2.8,-3.3) rectangle (-0.25,-2.8);
\draw[cadline] ( 0.25,-3.3) rectangle ( 2.8,-2.8);
\node[font=\tiny, right] at (3.1,-3.05) {\textbf{Pos.~5} Narbenplatte $E_2$};

%--- VERSCHLUSSKÖRPER (Pos. 1, oben) ---
\fill[hatch] (-3.0, 2.5) rectangle (-0.2, 5.5);
\fill[hatch] ( 0.2, 2.5) rectangle ( 3.0, 5.5);
\draw[cadthick] (-3.0, 2.5) rectangle (-0.2, 5.5);
\draw[cadthick] ( 0.2, 2.5) rectangle ( 3.0, 5.5);
% Deckel
\fill[hatch] (-3.0, 5.5) rectangle (3.0, 6.0);
\draw[cadthick] (-3.0, 5.5) rectangle (3.0, 6.0);
% Bohrung innen unten
\draw[cadthick] (-0.2, 1.8) -- (-0.2, 2.5);
\draw[cadthick] ( 0.2, 1.8) -- ( 0.2, 2.5);
% Narbenplatte E1 (Pos. 5, obere Ebene)
\fill[gray!60] (-2.8, 2.8) rectangle (-0.25, 3.3);
\fill[gray!60] ( 0.25, 2.8) rectangle ( 2.8, 3.3);
\draw[cadline] (-2.8, 2.8) rectangle (-0.25, 3.3);
\draw[cadline] ( 0.25, 2.8) rectangle ( 2.8, 3.3);
\node[font=\tiny, right] at (3.1, 3.05) {\textbf{Pos.~5} Narbenplatte $E_1$};

%--- PINS (Pos. 3, an Verschlusskörper) ---
% Stellvertretend 2 Pins links/rechts darstellen
\fill[pincolor!80] (-2.2, 2.3) rectangle (-1.8, 2.8);
\fill[pincolor!80] ( 1.8, 2.3) rectangle ( 2.2, 2.8);
\draw[cadline] (-2.2, 2.3) rectangle (-1.8, 2.8);
\draw[cadline] ( 1.8, 2.3) rectangle ( 2.2, 2.8);
\node[font=\tiny, left] at (-2.3, 2.55) {\textbf{Pos.~3} Pin};

%--- PIN-FEDER (Pos. 4) ---
\draw[cadline, decorate, decoration={coil, aspect=0.4, segment length=2pt, amplitude=2pt}]
  (-2.0, 1.5) -- (-2.0, 2.3);
\draw[cadline, decorate, decoration={coil, aspect=0.4, segment length=2pt, amplitude=2pt}]
  ( 2.0, 1.5) -- ( 2.0, 2.3);
\node[font=\tiny, left] at (-2.45, 1.9) {\textbf{Pos.~4}};

%--- ZENTRALDRUCKKNOPF (Pos. 6) ---
\fill[presscolor!50] (-0.5, 4.5) rectangle (0.5, 5.5);
\draw[cadthick] (-0.5, 4.5) rectangle (0.5, 5.5);
\node[font=\tiny, right] at (0.6, 5.0) {\textbf{Pos.~6}};

%--- KOPPELSTANGE (Pos. 7) ---
\fill[lockcolor!60] (-0.15, -1.8) rectangle (0.15, 4.5);
\draw[cadline] (-0.15, -1.8) rectangle (0.15, 4.5);
\node[font=\tiny, right] at (0.2, 1.3) {\textbf{Pos.~7} Koppelstange};

%--- RÃœCKSTELLFEDER (Pos. 8) ---
\draw[cadline, decorate, decoration={coil, aspect=0.4, segment length=3pt, amplitude=2.5pt}]
  (0, -2.5) -- (0, -1.8);
\node[font=\tiny, right] at (0.2, -2.15) {\textbf{Pos.~8}};

%--- O-RINGE (Pos. 10) ---
\fill[black!70] (-3.05, 0.0) ellipse (0.12 and 0.25);
\fill[black!70] ( 3.05, 0.0) ellipse (0.12 and 0.25);
\node[font=\tiny, left] at (-3.2, 0.0) {\textbf{Pos.~10}};

%--- SICHERUNGSRINGE (Pos. 9) ---
\fill[black!50] (-3.1, 2.3) rectangle (-2.9, 2.5);
\fill[black!50] ( 2.9, 2.3) rectangle ( 3.1, 2.5);
\node[font=\tiny, right] at (3.15, 2.4) {\textbf{Pos.~9}};

%--- MASSLINIE: Gesamthöhe ---
\draw[dim] (3.8, -6.0) -- (3.8, 6.0) node[midway, right, font=\tiny]{$H = 120{,}0$};
\draw[diml] (3.0, -6.0) -- (4.0, -6.0);
\draw[diml] (3.0,  6.0) -- (4.0,  6.0);

%--- MASSLINIE: Außendurchmesser ---
\draw[dim] (-3.0, 7.0) -- (3.0, 7.0) node[midway, above, font=\tiny]{$\varnothing\,D = 60{,}0$};
\draw[diml] (-3.0, 6.0) -- (-3.0, 7.2);
\draw[diml] ( 3.0, 6.0) -- ( 3.0, 7.2);

%--- MASSLINIE: Innendurchmesser ---
\draw[dim] (-0.2, -7.0) -- (0.2, -7.0) node[midway, below, font=\tiny]{$\varnothing\,d_i = 4{,}0$};
\draw[diml] (-0.2,-6.2) -- (-0.2,-7.2);
\draw[diml] ( 0.2,-6.2) -- ( 0.2,-7.2);

%--- MASSLINIE: Ebenenabstand ---
\draw[dim] (-4.2, -2.8) -- (-4.2, 2.8) node[midway, left, font=\tiny]{$\Delta z = 56{,}0$};
\draw[diml] (-3.0,-2.8) -- (-4.4,-2.8);
\draw[diml] (-3.0, 2.8) -- (-4.4, 2.8);

%--- MASSLINIE: Wandstärke ---
\draw[dim] (-3.0, -4.0) -- (-0.2, -4.0) node[midway, below, font=\tiny]{$t = 28{,}0$};
\draw[diml] (-3.0,-3.5) -- (-3.0,-4.2);
\draw[diml] (-0.2,-3.5) -- (-0.2,-4.2);

%--- TRENNLINIE zwischen Ober- und Unterring ---
\draw[cadthick, dashed, red!70] (-3.2, 0.0) -- (3.2, 0.0);
\node[red!70, font=\tiny\bfseries, right] at (3.3, 0.0) {Trennebene};

%--- Ebenen-Beschriftung ---
\node[blue!70, font=\tiny\bfseries, left] at (-3.2, 3.05) {$E_1$};
\node[green!60!black, font=\tiny\bfseries, left] at (-3.2,-3.05) {$E_2$};

%--- Titel und Rahmen ---
\node[font=\bfseries\small, above] at (0, 6.8) {SCHNITTANSICHT A--A \quad Maßstab 1:1};

\end{tikzpicture}
\caption{Längsschnitt durch die Mittelachse (Schnitt A--A). Dargestellt sind alle Hauptbauteile
         (Pos.~1--10) mit Hauptmaßen in mm. Schraffen: Vollmaterial V2A-Stahl. Grau: Narbenplatten.
         Rot gestrichelt: Trennebene zwischen Ober- und Unterring.}
\label{fig:cad_laengsschnitt}
\end{figure}

% -----------------------------------------------------------
\subsection{Draufsicht -- Verschlusskörper (Pos.~1, Oberring)}
% -----------------------------------------------------------

\begin{figure}[H]
\centering
\begin{tikzpicture}[scale=1.15,
  cadthick/.style={draw=black, line width=1.0pt},
  cadline/.style={draw=black, line width=0.4pt},
  cadcenter/.style={draw=black!40, line width=0.3pt, dash pattern=on 4pt off 2pt on 1pt off 2pt},
  cadthin/.style={draw=black!60, line width=0.25pt},
  dim/.style={<->, >=Stealth, thin},
  diml/.style={draw=black, thin},
  hatch/.style={pattern=north east lines, pattern color=black!35},
]

% Mittelachsen
\draw[cadcenter] (-3.5, 0) -- (3.5, 0);
\draw[cadcenter] (0, -3.5) -- (0, 3.5);

% Außenring
\draw[cadthick] (0,0) circle (3.0);
% Innenbohrung (Zentralkanal)
\draw[cadthick] (0,0) circle (0.5);
% Schraffur Vollmaterial (Ringquerschnitt, vereinfacht angedeutet)
\foreach \r in {2.6, 2.7, 2.8, 2.9} {
  \draw[cadthin, gray!30] (0,0) circle (\r);
}

% Zentraldruckknopf (Pos. 6)
\fill[presscolor!40] (0,0) circle (0.45);
\draw[cadthick, presscolor!80!black] (0,0) circle (0.45);
\node[font=\tiny, presscolor!80!black, below right] at (0.35, -0.35) {\textbf{6}};

% 9 Pins + Pinbohrungen auf Teilkreis r=2.1
\foreach \i in {1,...,9} {
  \pgfmathsetmacro{\ang}{(\i-1)*40}
  \pgfmathsetmacro{\rx}{2.1*cos(\ang)}
  \pgfmathsetmacro{\ry}{2.1*sin(\ang)}
  % Pinbohrung
  \fill[white] (\rx,\ry) circle (0.18);
  \draw[cadthick, pincolor] (\rx,\ry) circle (0.18);
  % Pin-Querschnitt (gefüllt)
  \fill[pincolor!70] (\rx,\ry) circle (0.13);
  \node[font=\tiny, black] at ({2.58*cos(\ang)},{2.58*sin(\ang)}) {$p_{\i}$};
}

% Teilkreis der Pins
\draw[cadcenter] (0,0) circle (2.1);
\node[cadcenter, font=\tiny, above right] at (1.6, 1.4) {$\varnothing\,T_p = 42{,}0$};

% 4 Befestigungsbohrungen (M4, Teilkreis r=2.55)
\foreach \ang in {25, 115, 205, 295} {
  \pgfmathsetmacro{\bx}{2.55*cos(\ang)}
  \pgfmathsetmacro{\by}{2.55*sin(\ang)}
  \fill[white] (\bx,\by) circle (0.10);
  \draw[cadline] (\bx,\by) circle (0.10);
  \draw[cadline] (\bx,\by) circle (0.16); % Senkung
}
\node[font=\tiny, gray] at (-2.5, -1.9) {$4\times$M4};

% Maßlinien
% Außendurchmesser
\draw[dim] (-3.0, -3.7) -- (3.0, -3.7) node[midway, below, font=\tiny]{$\varnothing\,60{,}0$};
\draw[diml] (-3.0,-3.0) -- (-3.0,-3.9);
\draw[diml] ( 3.0,-3.0) -- ( 3.0,-3.9);

% Innenbohrung
\draw[dim] (0.0, 0.0) -- (0.5, 0.0) node[midway, above, font=\tiny]{$r\,5{,}0$};

% Teilkreis Pins
\draw[dim, pincolor] (0,0) -- (2.1, 0) node[midway, below, font=\tiny\bfseries, pincolor]{$r_p = 21{,}0$};

% Winkelmaß Pin 1 zu Pin 2
\draw[dim, black!60] (0.65, 0.0) arc[start angle=0, end angle=40, radius=0.65]
  node[midway, right, font=\tiny]{$40°$};

% Titel
\node[font=\bfseries\small, above] at (0, 3.5) {DRAUFSICHT -- Verschlusskörper (Pos.~1) \quad M~1:1};

% Oberflächenangabe
\node[font=\tiny, right] at (3.2, 2.5) {$\sqrt{Ra\,0{,}8}$};
\node[font=\tiny, right] at (3.2, 2.1) {(alle Flächen)};

\end{tikzpicture}
\caption{Draufsicht auf den Verschlusskörper (Pos.~1, Oberring). Neun Pinbohrungen
         ($\varnothing\,3{,}0$ H7) auf Teilkreis $\varnothing\,42{,}0$\,mm, Winkelabstand $40°$.
         Zentralbohrung $\varnothing\,10{,}0$\,mm für Koppelstange. Vier Befestigungsbohrungen M4.}
\label{fig:cad_draufsicht_oberring}
\end{figure}

% -----------------------------------------------------------
\subsection{Narbenplatten-Detail -- L-förmige Rastgeometrie}
% -----------------------------------------------------------

\begin{figure}[H]
\centering
\begin{tikzpicture}[scale=2.5,
  cadthick/.style={draw=black, line width=0.8pt},
  cadline/.style={draw=black, line width=0.35pt},
  cadcenter/.style={draw=black!40, line width=0.25pt, dash pattern=on 3pt off 2pt on 1pt off 2pt},
  dim/.style={<->, >=Stealth, thin, black},
  diml/.style={draw=black, thin},
  hatch/.style={pattern=north east lines, pattern color=black!40},
  tol/.style={font=\tiny, black},
]

%--- L-FÖRMIGE NARBE: Einzeldarstellung (Maßstab 2:1) ---

% Materialkontur außen
\fill[scarcolor!15] (0,0) rectangle (2.2, 1.4);
\fill[white] (0.5,0) rectangle (0.9, 0.9);   % vertikaler Kanal (Einführ)
\fill[white] (0.5,0.6) rectangle (1.8, 0.9); % horizontaler Kanal (Rast)
\draw[cadthick, scarcolor!80!black]
  (0,0) -- (2.2,0) -- (2.2,1.4) -- (0,1.4) -- (0,0);
% Narbenkontur (L-Form Ausnehmung)
\draw[cadthick]
  (0.5,0) -- (0.5,0.6) -- (1.8,0.6) -- (1.8,0.9)
  -- (0.9,0.9) -- (0.9,0) -- cycle;

% Mittellinie Einführkanal
\draw[cadcenter] (0.7,-0.15) -- (0.7,1.0);
% Mittellinie Rastkanal
\draw[cadcenter] (0.2,0.75) -- (2.0,0.75);

% Radienmarkierung (Ecken-Radius R0.3)
\node[font=\tiny] at (0.53,0.63) {$R_{0.3}$};
\node[font=\tiny] at (0.87,0.63) {$R_{0.3}$};

% Maßlinien
% Breite Einführkanal
\draw[dim] (0.5,-0.22) -- (0.9,-0.22) node[midway, below, font=\tiny]{$b_E = 4{,}0$};
\draw[diml] (0.5,-0.05) -- (0.5,-0.28);
\draw[diml] (0.9,-0.05) -- (0.9,-0.28);

% Tiefe Einführkanal
\draw[dim] (-0.15,0) -- (-0.15,0.6) node[midway, left, font=\tiny]{$h_E = 6{,}0$};
\draw[diml] (0.0, 0.0) -- (-0.22, 0.0);
\draw[diml] (0.0, 0.6) -- (-0.22, 0.6);

% Breite Rastkanal
\draw[dim] (0.9,1.0) -- (1.8,1.0) node[midway, above, font=\tiny]{$b_R = 9{,}0$};
\draw[diml] (0.9,0.9) -- (0.9,1.08);
\draw[diml] (1.8,0.9) -- (1.8,1.08);

% Tiefe Rastkanal
\draw[dim] (2.1,0.6) -- (2.1,0.9) node[midway, right, font=\tiny]{$h_R = 3{,}0$};
\draw[diml] (1.85,0.6) -- (2.2,0.6);
\draw[diml] (1.85,0.9) -- (2.2,0.9);

% Gesamtbreite
\draw[dim] (0,-0.5) -- (2.2,-0.5) node[midway, below, font=\tiny]{$B = 22{,}0$};
\draw[diml] (0,-0.05) -- (0,-0.58);
\draw[diml] (2.2,-0.05) -- (2.2,-0.58);

% Gesamthöhe
\draw[dim] (2.45,0) -- (2.45,1.4) node[midway, right, font=\tiny]{$H_N = 14{,}0$};
\draw[diml] (2.2,0) -- (2.55,0);
\draw[diml] (2.2,1.4) -- (2.55,1.4);

% Tiefe Narbe (aus Blatt heraus, als Angabe)
\node[font=\tiny, right] at (2.6, 0.7) {Tiefe $t_N = 4{,}0$};
\node[font=\tiny, right] at (2.6, 0.5) {(senkr. zur Ebene)};

% Toleranzangaben
\node[tol, below] at (0.7, -0.65) {$b_E$: H7 / $+0{,}012$\,mm};
\node[tol, below] at (0.7, -0.82) {$b_R$: H8 / $+0{,}022$\,mm};

% Pin-Lage (gestrichelt)
\draw[cadline, blue!50, dashed] (0.7, 0.3) circle (0.2);
\node[font=\tiny, blue!60, right] at (0.92, 0.3) {Pin $\varnothing\,3{,}0$ (Einführlage)};

\draw[cadline, pincolor, dashed] (1.35, 0.75) circle (0.2);
\node[font=\tiny, pincolor, right] at (1.57, 0.75) {Pin $\varnothing\,3{,}0$ (Rastlage)};

% Maßstabshinweis
\node[font=\bfseries\small, above] at (1.1, 1.55) {NARBENDETAIL -- Pos.~5 \quad M~2:1};

% Drehrichtungs-Pfeil
\draw[->, thick, lockcolor!80!black] (0.7,0.3) -- (1.35, 0.73);
\node[lockcolor!80!black, font=\tiny] at (1.05, 0.65) {$\Phi_1$};

\end{tikzpicture}
\caption{Detailzeichnung der L-förmigen Verschluss-Narbe (Pos.~5, Maßstab 2:1). Einführkanal
         ($b_E \times h_E = 4{,}0 \times 6{,}0$\,mm) und Rastkanal ($b_R \times h_R = 9{,}0 \times 3{,}0$\,mm).
         Gestrichelt: Pin-Position bei Einführung (blau) und nach Einrastung (blau/dunkel).}
\label{fig:cad_narbe_detail}
\end{figure}

% -----------------------------------------------------------
\subsection{Pin-Zeichnung (Pos.~3) -- Einzelteil}
% -----------------------------------------------------------

\begin{figure}[H]
\centering
\begin{tikzpicture}[scale=2.2,
  cadthick/.style={draw=black, line width=0.8pt},
  cadline/.style={draw=black, line width=0.35pt},
  cadcenter/.style={draw=black!40, line width=0.25pt, dash pattern=on 3pt off 2pt on 1pt off 2pt},
  dim/.style={<->, >=Stealth, thin},
  diml/.style={draw=black, thin},
  hatch/.style={pattern=north east lines, pattern color=black!35},
]

%-- VORDERANSICHT: Zylindrischer Pin mit Rastnasen-Fase --
% Körper
\fill[pincolor!20] (-0.15, 0) rectangle (0.15, 1.4);
\draw[cadthick, pincolor!80!black] (-0.15, 0) rectangle (0.15, 1.4);
% Fase oben (Einführfase 45°)
\fill[pincolor!30] (-0.15,1.4) -- (0.15,1.4) -- (0.0,1.55) -- cycle;
\draw[cadthick] (-0.15,1.4) -- (0.0,1.55) -- (0.15,1.4);
% Rastnase (Hinterschneidung)
\fill[pincolor!50] (-0.15,0.55) -- (-0.25,0.55) -- (-0.25,0.65) -- (-0.15,0.75) -- cycle;
\fill[pincolor!50] ( 0.15,0.55) -- ( 0.25,0.55) -- ( 0.25,0.65) -- ( 0.15,0.75) -- cycle;
\draw[cadthick] (-0.15,0.55) -- (-0.25,0.55) -- (-0.25,0.65) -- (-0.15,0.75);
\draw[cadthick] ( 0.15,0.55) -- ( 0.25,0.55) -- ( 0.25,0.65) -- ( 0.15,0.75);
% Federsitz unten
\fill[pincolor!10] (-0.08,0) rectangle (0.08,-0.25);
\draw[cadline] (-0.08,0) rectangle (0.08,-0.25);
\node[font=\tiny, right] at (0.16,-0.12) {Federsitz};

% Mittelachse
\draw[cadcenter] (0,-0.4) -- (0,1.7);

% Maßlinien
% Gesamtlänge
\draw[dim] (0.5,0.0) -- (0.5,1.55) node[midway, right, font=\tiny]{$L = 15{,}5$};
\draw[diml] (0.15,0.0) -- (0.6,0.0);
\draw[diml] (0.15,1.55) -- (0.6,1.55);

% Durchmesser Körper
\draw[dim] (-0.15,-0.55) -- (0.15,-0.55) node[midway, below, font=\tiny]{$\varnothing\,d = 3{,}0\,\texttt{h6}$};
\draw[diml] (-0.15,-0.1) -- (-0.15,-0.62);
\draw[diml] ( 0.15,-0.1) -- ( 0.15,-0.62);

% Rastnase-Breite
\draw[dim] (-0.25,0.35) -- (0.25,0.35) node[midway, below, font=\tiny]{$d_R = 5{,}0$};
\draw[diml] (-0.25,0.55) -- (-0.25,0.28);
\draw[diml] ( 0.25,0.55) -- ( 0.25,0.28);

% Höhe Rastnase
\draw[dim] (-0.6,0.55) -- (-0.6,0.75) node[midway, left, font=\tiny]{$h_{\rm Rast} = 2{,}0$};
\draw[diml] (-0.25,0.55) -- (-0.68,0.55);
\draw[diml] (-0.25,0.75) -- (-0.68,0.75);

% Einführfase
\node[font=\tiny, right] at (0.27,1.48) {Fase $45° \times 0{,}3$};

% Werkstoff / Oberfl.
\node[font=\tiny, above] at (0,1.75) {\textbf{Pos.~3} -- Stift $\varnothing\,3{,}0 \times 15{,}5$\,mm};
\node[font=\tiny, above] at (0,1.62) {EN~1.4571 -- Härte: 240--280\,HV -- $Ra\,0{,}4$ (Lauffläche)};

%--- QUERSCHNITT A-A (rechts daneben) ---
\begin{scope}[xshift=1.3cm]
\fill[pincolor!20] (0,0) circle (0.15);
\draw[cadthick, pincolor!80!black] (0,0) circle (0.15);
\draw[cadcenter] (-0.35,0) -- (0.35,0);
\draw[cadcenter] (0,-0.35) -- (0,0.35);
\node[font=\tiny, below] at (0,-0.4) {Schnitt A--A};
\node[font=\tiny, below] at (0,-0.55) {$\varnothing\,3{,}0$};
\end{scope}

\node[font=\bfseries\small, above] at (0.5, 1.9) {EINZELTEIL -- Pin (Pos.~3) \quad M~4:1};

\end{tikzpicture}
\caption{Einzelteilzeichnung Pin (Pos.~3, Maßstab 4:1). Zylinderstift aus EN~1.4571 mit
         beidseitiger Rastnasen-Schulter ($d_R = 5{,}0$\,mm) zur Sicherung in der Narbe.
         Federsitz unten für Druckfeder (Pos.~4). Einführfase $45°$ oben.}
\label{fig:cad_pin}
\end{figure}

% -----------------------------------------------------------
\subsection{Zentraldruckknopf und Entriegelungsmechanik (Pos.~6/7)}
% -----------------------------------------------------------

\begin{figure}[H]
\centering
\begin{tikzpicture}[scale=1.3,
  cadthick/.style={draw=black, line width=0.8pt},
  cadline/.style={draw=black, line width=0.35pt},
  cadcenter/.style={draw=black!40, line width=0.25pt, dash pattern=on 3pt off 2pt on 1pt off 2pt},
  hatch/.style={pattern=north east lines, pattern color=black!40},
  dim/.style={<->, >=Stealth, thin},
  diml/.style={draw=black, thin},
]

% Mittelachse
\draw[cadcenter] (0,-5.0) -- (0,3.0);

%--- Zentraldruckknopf (Pos. 6) ---
\fill[presscolor!25] (-0.7,1.5) rectangle (0.7,2.5);
\fill[hatch] (-0.7,1.5) rectangle (0.7,2.5);
\draw[cadthick] (-0.7,1.5) rectangle (0.7,2.5);
% Fase
\fill[presscolor!30] (-0.7,2.5) -- (-0.5,2.8) -- (0.5,2.8) -- (0.7,2.5) -- cycle;
\draw[cadthick] (-0.7,2.5) -- (-0.5,2.8) -- (0.5,2.8) -- (0.7,2.5);
\node[font=\tiny, right] at (0.75, 2.1) {\textbf{Pos.~6} Druckknopf};

%--- Koppelstange (Pos. 7) ---
\fill[lockcolor!30] (-0.18,-4.0) rectangle (0.18,1.5);
\draw[cadthick] (-0.18,-4.0) rectangle (0.18,1.5);
\node[font=\tiny, right] at (0.22, -1.2) {\textbf{Pos.~7} Koppelstange};

% Querstift E1 (Aktivierung Narbenplatte E1)
\fill[pincolor!60] (-1.2, 0.5) rectangle (1.2, 0.8);
\draw[cadthick] (-1.2, 0.5) rectangle (1.2, 0.8);
\node[font=\tiny, left] at (-1.25, 0.65) {Querstift $E_1$};

% Querstift E2 (Aktivierung Narbenplatte E2)
\fill[pincolor!60] (-1.2, -3.2) rectangle (1.2, -2.9);
\draw[cadthick] (-1.2, -3.2) rectangle (1.2, -2.9);
\node[font=\tiny, left] at (-1.25, -3.05) {Querstift $E_2$};

%--- Rückstellfeder (Pos. 8) ---
\draw[cadline, decorate, decoration={coil, aspect=0.4, segment length=3.5pt, amplitude=3pt}]
  (0,-4.5) -- (0,-4.0);
\draw[cadline] (-0.25,-4.5) -- (0.25,-4.5);
\node[font=\tiny, right] at (0.25,-4.25) {\textbf{Pos.~8}};

% Narbenplatten-Betätiger (schematisch)
\fill[gray!40] (-1.5, 0.3) rectangle (-1.2, 1.0);
\fill[gray!40] ( 1.2, 0.3) rectangle ( 1.5, 1.0);
\draw[cadline] (-1.5, 0.3) rectangle (-1.2, 1.0);
\draw[cadline] ( 1.2, 0.3) rectangle ( 1.5, 1.0);
\fill[gray!40] (-1.5,-3.4) rectangle (-1.2,-2.7);
\fill[gray!40] ( 1.2,-3.4) rectangle ( 1.5,-2.7);
\draw[cadline] (-1.5,-3.4) rectangle (-1.2,-2.7);
\draw[cadline] ( 1.2,-3.4) rectangle ( 1.5,-2.7);

% Ebenen-Beschriftung
\draw[blue!60, dashed, thick] (-1.7, 0.65) -- (1.7, 0.65);
\node[blue!70, font=\tiny\bfseries, left] at (-1.75, 0.65) {$E_1$};
\draw[green!60!black, dashed, thick] (-1.7,-3.05) -- (1.7,-3.05);
\node[green!60!black, font=\tiny\bfseries, left] at (-1.75,-3.05) {$E_2$};

% Druckpfeil
\draw[->, ultra thick, red!70, line width=2pt] (0, 3.6) -- (0, 2.85);
\node[red!70, font=\bfseries\small, right] at (0.1, 3.3) {$F_\pi$};

% Maßlinien
\draw[dim] (2.0, 0.65) -- (2.0,-3.05) node[midway, right, font=\tiny]{$\Delta z = 56{,}0$};
\draw[diml] (1.7, 0.65) -- (2.2, 0.65);
\draw[diml] (1.7,-3.05) -- (2.2,-3.05);

\draw[dim] (-2.2,-4.5) -- (-2.2, 2.8) node[midway, left, font=\tiny]{$L_K = 90{,}0$};
\draw[diml] (-0.18,-4.5) -- (-2.4,-4.5);
\draw[diml] (-0.5, 2.8) -- (-2.4, 2.8);

\node[font=\bfseries\small, above] at (0, 3.8) {ENTRIEGELUNGSMECHANIK -- Pos.~6/7 \quad M~1:1};
\node[font=\tiny, below] at (0,-5.0) {Betätigung: $F_\pi$ > $F_{\rm Rast,ges}$ $\Rightarrow$ simultane Freigabe $E_1$ und $E_2$};

\end{tikzpicture}
\caption{Schnittdarstellung der Entriegelungsmechanik (Pos.~6 Druckknopf, Pos.~7 Koppelstange).
         Druckimpuls $F_\pi$ überträgt sich über die Koppelstange simultan auf beide
         Narbenplatten-Betätiger ($E_1$ und $E_2$). Rückstellfeder Pos.~8 gewährleistet Rückhub.}
\label{fig:cad_entriegelung}
\end{figure}

% -----------------------------------------------------------
\subsection{Toleranz- und Passungskonzept}
% -----------------------------------------------------------
Tabelle \ref{tab:toleranzen} zeigt die Toleranz- und Passungstabelle für alle kritischen Fügestellen des DRLM.
\begin{sidewaystable}
  \centering
  \caption{Toleranz- und Passungstabelle für alle kritischen Fügestellen des DRLM.}
  \label{tab:toleranzen}
  \renewcommand{\arraystretch}{1.35}
  \begin{tabular}{llllll}
    \toprule
    \textbf{Fügestelle} & \textbf{Welle/Bohrung} & \textbf{Nennmaß [mm]} &
    \textbf{Passung} & \textbf{Spiel/Übermaß [$\mu$m]} & \textbf{Funktion} \\
    \midrule
    Pin -- Pinbohrung         & $\varnothing\,3{,}0$ & 3{,}0  & H7/h6 & $+0$\ldots$+22$   & Gleitpassung \\
    Pin -- Einführkanal       & $b_E = 4{,}0$        & 4{,}0  & H8/f7 & $+16$\ldots$+50$  & Leichtlauf \\
    Pin -- Rastkanal          & $b_R \times h_R$     & 3{,}0  & H7/h6 & $+2$\ldots$+20$   & Rastpassung \\
    Koppelstange -- Führung   & $\varnothing\,3{,}6$ & 3{,}6  & H7/g6 & $+2$\ldots$+26$   & Schiebepassung \\
    Druckknopf -- Gehäuse     & $\varnothing\,10{,}0$& 10{,}0 & H8/f7 & $+13$\ldots$+49$  & Leichtlauf \\
    Oberring -- Unterring     & $\varnothing\,60{,}0$& 60{,}0 & H7/g6 & $+10$\ldots$+39$  & Drehlaufpassung \\
    Sicherungsring -- Nut     & $\varnothing\,42{,}0$& 42{,}0 & ---   & nach DIN~471       & Axialsicherung \\
    O-Ring -- Nut             & Breite $3{,}6$       & 3{,}6  & ---   & nach DIN~3771      & Abdichtung \\
    \bottomrule
  \end{tabular}
\end{sidewaystable}

% -----------------------------------------------------------
\subsection{Oberflächenangaben und Fertigungshinweise}
% -----------------------------------------------------------

\begin{table}[H]
  \centering
  \caption{Oberflächenangaben und Fertigungshinweise für alle Bauteile.}
  \label{tab:oberflaechen}
  \renewcommand{\arraystretch}{1.35}
  \begin{tabular}{llll}
    \toprule
    \textbf{Bauteil (Pos.)} & \textbf{Fläche} & \textbf{Ra [$\mu$m]} & \textbf{Fertigungshinweis} \\
    \midrule
    Verschlusskörper (1)  & Lauffläche (innen, $\varnothing$60)  & 0{,}4 & Drehen, Schleifen \\
    Verschlusskörper (1)  & Pinbohrungen ($\varnothing$3 H7)     & 0{,}8 & Reiben \\
    Verschlusskörper (1)  & Deckfläche                           & 1{,}6 & Drehen \\
    Sockelring (2)        & Lauffläche (außen, $\varnothing$60)  & 0{,}4 & Drehen, Schleifen \\
    Sockelring (2)        & Narbenkanäle                         & 0{,}8 & Fräsen, Senkerodieren \\
    Pin (3)               & Zylinderfläche ($\varnothing$3 h6)   & 0{,}4 & Schleifen \\
    Pin (3)               & Rastnasen-Schulter                   & 0{,}8 & Fräsen \\
    Narbenplatte (5)      & Kanalflächen                         & 0{,}8 & Fräsen \\
    Druckknopf (6)        & Außenfläche                          & 0{,}8 & Drehen \\
    Koppelstange (7)      & Führungsfläche                       & 0{,}4 & Schleifen \\
    \bottomrule
  \end{tabular}
\end{table}

\subsubsection*{Allgemeine Fertigungshinweise}

\begin{enumerate}[noitemsep]
  \item \textbf{Werkstoff}: Alle tragenden Teile aus EN~1.4301 (AISI~304) oder EN~1.4571
        (AISI~316Ti); Federn aus Federstahl 1.4310.
  \item \textbf{Wärmebehandlung}: Pins (Pos.~3) nach Fertigung Lösungsglühen bei $1050\,°C$
        + Abschrecken (Wasser), Zielzustand: $240$--$280\,\text{HV}$.
  \item \textbf{Allgemeintoleranzen}: ISO~2768-m (mittlere Toleranzklasse) für alle
        nicht bemaßten Abmessungen.
  \item \textbf{Formtoleranzen}: Rundheit der Pinbohrungen $\leq 0{,}005\,\text{mm}$;
        Koaxialität Koppelstange zu Mittelachse $\leq 0{,}02\,\text{mm}$.
  \item \textbf{Montagereihenfolge}: (a)~Narbenplatten in Sockelring einsetzen,
        (b)~Federn und Pins einsetzen, (c)~Koppelstange einführen,
        (d)~O-Ringe einlegen, (e)~Oberring aufsetzen und Sicherungsringe montieren.
  \item \textbf{Funktionsprüfung}: Rastmoment bei $\Phi_1$: $0{,}3$--$0{,}6\,\text{Nm}$;
        Entriegelungskraft $F_\pi$: $25$--$45\,\text{N}$; Haltekraft $F_{\rm Rast,ges}$: $\geq 500\,\text{N}$.
  \item \textbf{Korrosionsschutz}: Alle Teile nach Montage passivieren (HNO\textsubscript{3}-Bad,
        30~min, $50\,°C$) gemäß ASTM~A~967.
\end{enumerate}

% -----------------------------------------------------------
\subsection{Montagezeichnung -- Explosionsdarstellung}
% -----------------------------------------------------------

\begin{figure}[H]
\centering
\begin{tikzpicture}[scale=0.95,
  cadthick/.style={draw=black, line width=0.8pt},
  cadline/.style={draw=black, line width=0.35pt},
  montage/.style={draw=gray!60, dashed, line width=0.4pt},
  hatch/.style={pattern=north east lines, pattern color=black!35},
]

% --- EXPLOSIONSDARSTELLUNG (schematisch, isometrisch angedeutet) ---

%---- Sockelring (Pos. 2) -- unten ----
\fill[gray!20, draw=black, cadthick] (-3.0,-6.0) rectangle (3.0,-3.5);
\fill[white] (-0.5,-6.0) rectangle (0.5,-3.5); % Innenkanal
\draw[cadthick] (-0.5,-6.0) rectangle (0.5,-3.5);
\node[font=\small\bfseries] at (-4.0,-4.75) {Pos.~2};
\node[font=\small] at (-4.0,-5.15) {Sockelring};

%---- Narbenplatte E2 (Pos. 5) ----
\fill[scarcolor!20, draw=scarcolor, cadthick] (-2.8,-3.2) rectangle (2.8,-2.7);
\fill[white] (-2.0,-3.2) rectangle (-1.5,-2.7);
\fill[white] (-0.5,-3.2) rectangle (0.5,-2.7);
\fill[white] (1.5,-3.2) rectangle (2.0,-2.7);
\node[font=\small\bfseries, scarcolor!80!black] at (4.0,-2.95) {Pos.~5};
\node[font=\small, scarcolor!80!black] at (4.0,-3.3) {Narbenplatte $E_2$};
\draw[montage] (3.0,-2.95) -- (3.3,-2.95);

%---- Rückstellfeder (Pos. 8) ----
\draw[cadline, decorate, decoration={coil, aspect=0.45, segment length=4pt, amplitude=3.5pt}]
  (0,-2.4) -- (0,-1.5);
\node[font=\small\bfseries] at (1.5,-1.95) {Pos.~8};
\draw[montage] (0.3,-1.95) -- (1.2,-1.95);

%---- Koppelstange (Pos. 7) ----
\fill[lockcolor!40, draw=black, cadthick] (-0.25,-1.4) rectangle (0.25,2.3);
\node[font=\small\bfseries] at (1.5,0.5) {Pos.~7};
\draw[montage] (0.3,0.5) -- (1.2,0.5);

%---- Narbenplatte E1 (Pos. 5) ----
\fill[scarcolor!20, draw=scarcolor, cadthick] (-2.8,2.4) rectangle (2.8,2.9);
\fill[white] (-2.0,2.4) rectangle (-1.5,2.9);
\fill[white] (-0.5,2.4) rectangle (0.5,2.9);
\fill[white] (1.5,2.4) rectangle (2.0,2.9);
\node[font=\small\bfseries, scarcolor!80!black] at (4.0,2.65) {Pos.~5};
\node[font=\small, scarcolor!80!black] at (4.0,2.3) {Narbenplatte $E_1$};
\draw[montage] (3.0,2.65) -- (3.3,2.65);

%---- Pins + Federn (3 sichtbar, stellvertretend) ----
\fill[pincolor!70, draw=black, cadthick] (-2.1,2.9) rectangle (-1.9,4.1);
\fill[pincolor!70, draw=black, cadthick] (-0.1,2.9) rectangle ( 0.1,4.1);
\fill[pincolor!70, draw=black, cadthick] ( 1.9,2.9) rectangle ( 2.1,4.1);
\draw[cadline, decorate, decoration={coil, aspect=0.45, segment length=2.5pt, amplitude=2pt}]
  (-2.0,3.5) -- (-2.0,2.9);
\draw[cadline, decorate, decoration={coil, aspect=0.45, segment length=2.5pt, amplitude=2pt}]
  (0.0, 3.5) -- (0.0, 2.9);
\draw[cadline, decorate, decoration={coil, aspect=0.45, segment length=2.5pt, amplitude=2pt}]
  (2.0, 3.5) -- (2.0, 2.9);
\node[font=\small\bfseries, pincolor!80!black] at (-4.0, 3.5) {Pos.~3+4};
\node[font=\small, pincolor!80!black] at (-4.0, 3.1) {Pins + Federn};

%---- Verschlusskörper / Oberring (Pos. 1) ----
\fill[gray!20, draw=black, cadthick] (-3.0,4.2) rectangle (3.0,6.0);
\fill[white] (-0.5,4.2) rectangle (0.5,6.0);
\draw[cadthick] (-0.5,4.2) rectangle (0.5,6.0);
\node[font=\small\bfseries] at (-4.0,5.1) {Pos.~1};
\node[font=\small] at (-4.0,4.7) {Verschlusskörper};

%---- Zentraldruckknopf (Pos. 6) ----
\fill[presscolor!40, draw=black, cadthick] (-0.6,6.1) rectangle (0.6,7.1);
\draw[cadthick] (-0.6,7.1) -- (-0.4,7.5) -- (0.4,7.5) -- (0.6,7.1);
\node[font=\small\bfseries, presscolor!80!black] at (1.8, 6.8) {Pos.~6};
\node[font=\small, presscolor!80!black] at (1.8, 6.4) {Druckknopf};

%---- Sicherungsringe (Pos. 9, angedeutet) ----
\fill[black!50, draw=black] (-3.1,4.0) rectangle (-2.9,4.2);
\fill[black!50, draw=black] ( 2.9,4.0) rectangle ( 3.1,4.2);
\node[font=\tiny] at (-4.0,4.05) {Pos.~9 Sicherungsring};

%---- O-Ringe (Pos. 10) ----
\fill[black!60] (-3.15,-1.5) ellipse (0.12 and 0.3);
\fill[black!60] ( 3.15,-1.5) ellipse (0.12 and 0.3);
\node[font=\tiny] at (4.5,-1.5) {Pos.~10 O-Ring};

%---- Montagepfeile (gestrichelt) ----
\draw[->, montage, thick] (0, 4.0) -- (0, 3.1);
\draw[->, montage, thick] (0, 2.3) -- (0, 1.2);
\draw[->, montage, thick] (0,-1.4) -- (0,-2.3);
\draw[->, montage, thick] (0,-2.5) -- (0,-3.1);

% Rahmen
\node[font=\bfseries\small, above] at (0, 7.8)
  {EXPLOSIONSANSICHT -- DRLM Gesamtmontage \quad M~1:1};
\node[font=\tiny, below] at (0,-6.5)
  {Montagereihenfolge: $2 \to 5_{\rm unten} \to 8 \to 7 \to 5_{\rm oben} \to 3{+}4 \to 1 \to 6 \to 9 \to 10$};

\end{tikzpicture}
\caption{Explosionsdarstellung des DRLM (schematischer Längsschnitt). Die Montagepfeile
         zeigen die empfohlene Montagereihenfolge. Grau: Stahlteile; Rot: Narbenplatten;
         Blau: Pins; Gelb: Koppelstange; Violett: Druckknopf.}
\label{fig:cad_explosion}
\end{figure}

% -----------------------------------------------------------
\subsection{Schriftfeld und Normangaben (Zeichnungsrahmen)}
% -----------------------------------------------------------

\begin{figure}[H]
\centering
\begin{tikzpicture}[scale=1.0, font=\tiny]
  % Äußerer Rahmen
  \draw[line width=1.5pt] (0,0) rectangle (16, 5.5);
  % Innere Trennlinien
  \draw[thin] (10,0) -- (10,5.5);
  \draw[thin] (0,3.5) -- (10,3.5);
  \draw[thin] (0,2.0) -- (10,2.0);
  \draw[thin] (0,0.8) -- (16,0.8);
  \draw[thin] (10,2.0) -- (10,5.5);
  \draw[thin] (13,0.8) -- (13,5.5);
  % Vertikale Linien Schriftfeld rechts
  \draw[thin] (10,4.5) -- (16,4.5);
  \draw[thin] (10,3.2) -- (16,3.2);
  \draw[thin] (10,1.8) -- (16,1.8);

  % Felder Beschriftung
  \node[align=left, anchor=north west] at (0.15, 5.4)
    {\textbf{Benennung:}};
  \node[align=left, anchor=north west] at (0.15, 4.8)
    {\LARGE\bfseries Doppelter Ringverschluss};
  \node[align=left, anchor=north west] at (0.15, 3.9)
    {\textbf{Norm:} DIN~ISO~128 \quad \textbf{Maßstab:} 1:1 / 2:1 / 4:1};
  \node[align=left, anchor=north west] at (0.15, 3.3)
    {\textbf{Werkstoff (Hauptteile):} EN~1.4301 (V2A-Stahl)};
  \node[align=left, anchor=north west] at (0.15, 2.7)
    {\textbf{Allgemeintoleranzen:} ISO~2768-m};
  \node[align=left, anchor=north west] at (0.15, 2.2)
    {\textbf{Oberflächengüte (allg.):} $Ra\,1{,}6$};

  % Projektionsmethode
  \node[align=left, anchor=north west] at (0.15, 1.7)
    {\textbf{Projektionsmethode:} Europäisch (1.~Winkel) \quad \textbf{Einheit:} mm};

  % Schriftfeld rechts oben
  \node[align=left, anchor=north west] at (10.15, 5.4)
    {\textbf{Unternehmen / Institution:}};
  \node[align=left, anchor=north west] at (10.15, 4.8)
    {Technische Analyse DRLM};
  \node[align=left, anchor=north west] at (10.15, 4.4)
    {\textbf{Zeichnungs-Nr.:} DRG-000};

  \node[align=left, anchor=north west] at (10.15, 3.7)
    {\textbf{Erstellt:} --- \quad \textbf{Geprüft:} ---};
  \node[align=left, anchor=north west] at (10.15, 3.4)
    {\textbf{Datum:} \today};

  \node[align=left, anchor=north west] at (10.15, 2.8)
    {\textbf{Blatt:} 1 von 6 \quad \textbf{Rev.:} A};
  \node[align=left, anchor=north west] at (10.15, 2.2)
    {\textbf{Format:} A4};

  % Unten
  \node[align=left, anchor=north west] at (0.15, 0.75)
    {\textbf{Änderungen:} -- \quad (Rev.~A: Erstausgabe)};
  \node[align=right, anchor=north east] at (15.9, 0.75)
    {DRG-000 / Rev.~A};

  % Projektionssymbol (vereinfacht)
  \begin{scope}[xshift=12.5cm, yshift=2.5cm, scale=0.35]
    \draw[thin] (0,0) -- (1,0.5) -- (1,-0.5) -- cycle;
    \draw[thin] (1.5,0) circle (0.5);
    \node[below, font=\tiny] at (0.75,-0.8) {1. Winkel};
  \end{scope}
\end{tikzpicture}
\caption{Normgerechtes Schriftfeld (DIN~ISO~7200) für den Zeichnungssatz DRLM.
         Zeichnungsnummern DRG-001 bis DRG-006 für Einzelteilzeichnungen; DRG-000 für Gesamtmontage.}
\label{fig:schriftfeld}
\end{figure}

\end{document}


